\chapter{Introduction}
\section{Background}
Digital forensic examination depends on understanding what a storage format
makes visible and what a particular tool chooses to display. A normal file
listing is useful for navigation, but it is not a complete representation of
the bytes associated with an NTFS file. Microsoft documents files as having
one default stream and potentially additional named streams \cite{streams,clusters}.
For this project, the distinction matters because an examiner who exports
only the default stream can recover a valid-looking host file while omitting
the controlled information attached to it.

Information hiding also has a longer history than digital filesystems.
Anderson describes historical concealed messages and distinguishes passive
observers from opponents able to modify communications \cite{anderson}.
The relevant lesson for this laboratory is methodological: successfully
storing a hidden message does not establish that it will remain secret,
survive alteration, or be recoverable under every acquisition method. Each
claim requires its own experiment. Our work addresses recovery and integrity
in a deliberately limited storage environment.

The second location examined here is final-cluster file slack: bytes after
logical content ends in its final allocated cluster. We established a known
marker's presence through recovery and comparison with its earlier recorded
hash, rather than inferring it solely from unused capacity.

\section{Problem Statement}
The project asks whether a purpose-built Python parser can interpret the
necessary structures in a controlled raw NTFS image, locate documented ADS
and slack payloads, and recover them without changing the source image.
A related problem is how to report the strength of the evidence: successful
file identification, matching length, matching original hash, and unchanged
source image support different conclusions and must not be merged into one
unqualified success claim.

The contribution is an inspectable educational workflow connecting on-disk
offsets, extraction decisions, and verification reports. It does not claim
that mature tools lack ADS or slack capabilities.

\section{Aims and Objectives}
The overall aim is to demonstrate controlled NTFS information hiding and
defensive recovery with evidence that another reader can inspect. Six
measurable objectives guide the evaluation:
\begin{enumerate}
\item To characterize the laboratory image by recording its hash, partition
start, sector and cluster sizes, MFT location, and record size.
\item To identify the eight known baseline files and compare their names and
logical lengths with independently recorded Windows values.
\item To document a named ADS marker and a bounded slack marker, including
their lengths, locations, and expected SHA-256 values.
\item To recover both hidden markers exactly while retaining the original
source image and writing outputs to separate new files.
\item To test malformed inputs, unsupported layouts, output protection, and
verification outcomes using reproducible automated cases.
\item To package the source, dataset-generation scripts, instructions, and
results so a fresh checkout can reproduce the supported workflow.
\end{enumerate}

Original/recovered SHA-256 matches for all eight ordinary files were
confirmed by the team on November 25. These comparisons remain separate
from the independently documented hidden-marker comparisons.

\section{Project Scope and Limitations}
The primary experiment uses a dedicated VirtualBox NTFS disk exported as a
raw image. The implementation supports selected in-use MFT records,
resident data, ordinary nonresident extents, exact named-stream selection,
and bounded reads from the final cluster's slack. It does not reconstruct
deleted files, decrypt BitLocker or EFS, resolve attribute lists, or traverse
arbitrary fragmented MFT storage. A successful controlled experiment is not
evidence of compatibility with every NTFS volume.

Handcrafted NTFS-shaped fixtures exercise supported paths without the
large image; they are not mountable Windows filesystems. The 50 synthetic
payload samples vary bytes and lengths within one layout, rather than
representing independently acquired volumes.

Two fixed detection rules were evaluated separately on 100 labeled synthetic
images: named-stream screening and a known-marker search in calculated slack.
Their measured false-positive and evasion rates apply only to the constructed
cases and selected records. No statistical steganalysis classifier,
confidentiality layer, or general unknown-payload detector was implemented.
Chapter 5 explains the measured storage and detection metrics and why a
visual-quality score is not applicable to this filesystem carrier.

\section{Thesis Outline}
Chapter 2 reviews information hiding, NTFS structures, and relevant tools.
Chapter 3 defines the architecture, extraction rules, formulas, and
experimental design. Chapter 4 explains the code and development challenges.
Chapter 5 presents historical laboratory results and the new synthetic
campaign. Chapter 6 examines attacks and limitations. Chapter 7 evaluates
the objectives and proposes future work. The appendices provide execution
instructions, raw results, and the development record.

\section{Responsible Disclosure and Ethical Use Statement}
The software and experiments are intended solely for education, research,
and defensive analysis of authorized data. All hiding operations were
restricted to dedicated laboratory material or newly generated fixtures.
The tools must not be used for malicious concealment, unauthorized access,
or alteration of another person's evidence. Recovery should operate on
preserved images, with separate destinations for generated files.

NIST places forensic activity within an organization's documented process
\cite{nist86}. Accordingly, image identities and evidence gaps are retained.
This project is not certified for operational casework. No third-party
vulnerability requiring disclosure was discovered.
